\documentclass[]{article}
\usepackage{lmodern}
\usepackage{amssymb,amsmath}
\usepackage{ifxetex,ifluatex}
\usepackage{fixltx2e} % provides \textsubscript
\ifnum 0\ifxetex 1\fi\ifluatex 1\fi=0 % if pdftex
  \usepackage[T1]{fontenc}
  \usepackage[utf8]{inputenc}
\else % if luatex or xelatex
  \ifxetex
    \usepackage{mathspec}
  \else
    \usepackage{fontspec}
  \fi
  \defaultfontfeatures{Ligatures=TeX,Scale=MatchLowercase}
\fi
% use upquote if available, for straight quotes in verbatim environments
\IfFileExists{upquote.sty}{\usepackage{upquote}}{}
% use microtype if available
\IfFileExists{microtype.sty}{%
\usepackage{microtype}
\UseMicrotypeSet[protrusion]{basicmath} % disable protrusion for tt fonts
}{}
\usepackage[left=1.1in, right=1.1in, top=1.1in, bottom=1.1in]{geometry}
\usepackage{hyperref}
\hypersetup{unicode=true,
            pdftitle={Stat 435 Lecture Notes 4},
            pdfauthor={Xiongzhi Chen; Washington State University},
            pdfborder={0 0 0},
            breaklinks=true}
\urlstyle{same}  % don't use monospace font for urls
\usepackage{color}
\usepackage{fancyvrb}
\newcommand{\VerbBar}{|}
\newcommand{\VERB}{\Verb[commandchars=\\\{\}]}
\DefineVerbatimEnvironment{Highlighting}{Verbatim}{commandchars=\\\{\}}
% Add ',fontsize=\small' for more characters per line
\usepackage{framed}
\definecolor{shadecolor}{RGB}{248,248,248}
\newenvironment{Shaded}{\begin{snugshade}}{\end{snugshade}}
\newcommand{\KeywordTok}[1]{\textcolor[rgb]{0.13,0.29,0.53}{\textbf{{#1}}}}
\newcommand{\DataTypeTok}[1]{\textcolor[rgb]{0.13,0.29,0.53}{{#1}}}
\newcommand{\DecValTok}[1]{\textcolor[rgb]{0.00,0.00,0.81}{{#1}}}
\newcommand{\BaseNTok}[1]{\textcolor[rgb]{0.00,0.00,0.81}{{#1}}}
\newcommand{\FloatTok}[1]{\textcolor[rgb]{0.00,0.00,0.81}{{#1}}}
\newcommand{\ConstantTok}[1]{\textcolor[rgb]{0.00,0.00,0.00}{{#1}}}
\newcommand{\CharTok}[1]{\textcolor[rgb]{0.31,0.60,0.02}{{#1}}}
\newcommand{\SpecialCharTok}[1]{\textcolor[rgb]{0.00,0.00,0.00}{{#1}}}
\newcommand{\StringTok}[1]{\textcolor[rgb]{0.31,0.60,0.02}{{#1}}}
\newcommand{\VerbatimStringTok}[1]{\textcolor[rgb]{0.31,0.60,0.02}{{#1}}}
\newcommand{\SpecialStringTok}[1]{\textcolor[rgb]{0.31,0.60,0.02}{{#1}}}
\newcommand{\ImportTok}[1]{{#1}}
\newcommand{\CommentTok}[1]{\textcolor[rgb]{0.56,0.35,0.01}{\textit{{#1}}}}
\newcommand{\DocumentationTok}[1]{\textcolor[rgb]{0.56,0.35,0.01}{\textbf{\textit{{#1}}}}}
\newcommand{\AnnotationTok}[1]{\textcolor[rgb]{0.56,0.35,0.01}{\textbf{\textit{{#1}}}}}
\newcommand{\CommentVarTok}[1]{\textcolor[rgb]{0.56,0.35,0.01}{\textbf{\textit{{#1}}}}}
\newcommand{\OtherTok}[1]{\textcolor[rgb]{0.56,0.35,0.01}{{#1}}}
\newcommand{\FunctionTok}[1]{\textcolor[rgb]{0.00,0.00,0.00}{{#1}}}
\newcommand{\VariableTok}[1]{\textcolor[rgb]{0.00,0.00,0.00}{{#1}}}
\newcommand{\ControlFlowTok}[1]{\textcolor[rgb]{0.13,0.29,0.53}{\textbf{{#1}}}}
\newcommand{\OperatorTok}[1]{\textcolor[rgb]{0.81,0.36,0.00}{\textbf{{#1}}}}
\newcommand{\BuiltInTok}[1]{{#1}}
\newcommand{\ExtensionTok}[1]{{#1}}
\newcommand{\PreprocessorTok}[1]{\textcolor[rgb]{0.56,0.35,0.01}{\textit{{#1}}}}
\newcommand{\AttributeTok}[1]{\textcolor[rgb]{0.77,0.63,0.00}{{#1}}}
\newcommand{\RegionMarkerTok}[1]{{#1}}
\newcommand{\InformationTok}[1]{\textcolor[rgb]{0.56,0.35,0.01}{\textbf{\textit{{#1}}}}}
\newcommand{\WarningTok}[1]{\textcolor[rgb]{0.56,0.35,0.01}{\textbf{\textit{{#1}}}}}
\newcommand{\AlertTok}[1]{\textcolor[rgb]{0.94,0.16,0.16}{{#1}}}
\newcommand{\ErrorTok}[1]{\textcolor[rgb]{0.64,0.00,0.00}{\textbf{{#1}}}}
\newcommand{\NormalTok}[1]{{#1}}
\usepackage{graphicx,grffile}
\makeatletter
\def\maxwidth{\ifdim\Gin@nat@width>\linewidth\linewidth\else\Gin@nat@width\fi}
\def\maxheight{\ifdim\Gin@nat@height>\textheight\textheight\else\Gin@nat@height\fi}
\makeatother
% Scale images if necessary, so that they will not overflow the page
% margins by default, and it is still possible to overwrite the defaults
% using explicit options in \includegraphics[width, height, ...]{}
\setkeys{Gin}{width=\maxwidth,height=\maxheight,keepaspectratio}
\IfFileExists{parskip.sty}{%
\usepackage{parskip}
}{% else
\setlength{\parindent}{0pt}
\setlength{\parskip}{6pt plus 2pt minus 1pt}
}
\setlength{\emergencystretch}{3em}  % prevent overfull lines
\providecommand{\tightlist}{%
  \setlength{\itemsep}{0pt}\setlength{\parskip}{0pt}}
\setcounter{secnumdepth}{0}
% Redefines (sub)paragraphs to behave more like sections
\ifx\paragraph\undefined\else
\let\oldparagraph\paragraph
\renewcommand{\paragraph}[1]{\oldparagraph{#1}\mbox{}}
\fi
\ifx\subparagraph\undefined\else
\let\oldsubparagraph\subparagraph
\renewcommand{\subparagraph}[1]{\oldsubparagraph{#1}\mbox{}}
\fi

%%% Use protect on footnotes to avoid problems with footnotes in titles
\let\rmarkdownfootnote\footnote%
\def\footnote{\protect\rmarkdownfootnote}

%%% Change title format to be more compact
\usepackage{titling}

% Create subtitle command for use in maketitle
\newcommand{\subtitle}[1]{
  \posttitle{
    \begin{center}\large#1\end{center}
    }
}

\setlength{\droptitle}{-2em}

  \title{Stat 435 Lecture Notes 4}
    \pretitle{\vspace{\droptitle}\centering\huge}
  \posttitle{\par}
    \author{Xiongzhi Chen \\ Washington State University}
    \preauthor{\centering\large\emph}
  \postauthor{\par}
    \date{}
    \predate{}\postdate{}
  
\usepackage{bbm}
\usepackage{amssymb}
\usepackage{amsmath}
\usepackage{graphicx,float}

\begin{document}
\maketitle

{
\setcounter{tocdepth}{2}
\tableofcontents
}
\section{}\label{section}

\section{Bootstrap: motivation}\label{bootstrap-motivation}

\subsection{Overview}\label{overview}

\begin{itemize}
\item
  The \emph{bootsrtap} is mainly used to \emph{estimate and quantify the
  uncertainty associated with a given estimate or statistical learning
  methed}
\item
  For example, it can be used to estimate the standard error of an
  estimate (such as an estimated coefficient in a regression model)
\item
  The bootstrap \emph{may not work well when sample size is small or
  when sample comes from a relatively small region of the distribution
  of an unknown data generating process}
\end{itemize}

\subsection{Illustration I: problem}\label{illustration-i-problem}

Problem formulation:

\begin{itemize}
\tightlist
\item
  suppose we wish to invest a fixed sum of money into two financial
  assets that yield (random) returns of \(X\) and \(Y\), respectively
\item
  we will invest a fraction \(\alpha\) of our money in \(X\), and the
  rest \(1-\alpha\) in \(Y\)
\item
  we need to choose \(\alpha\) that minimizes the risk, or variance, of
  our investment.
\end{itemize}

Namely, we need to find \(\alpha\) that minimizes
\[\textrm{Var}(\alpha X + (1-\alpha)Y)\]

\subsection{Illustration I: solution}\label{illustration-i-solution}

\begin{itemize}
\item
  By calculus, we know that \[
  \alpha=\frac{\sigma_{Y}^{2}-\sigma_{XY}}{\sigma_{X}^{2}+\sigma_{Y}^{2}-2\sigma_{XY}}
  \] minimizes \[\textrm{Var}(\alpha X + (1-\alpha)Y),\] where
  \(\sigma_{X}^{2}=\textrm{Var}(X)\), \(\sigma_{Y}^{2}=\textrm{Var}(Y)\)
  and \(\sigma_{XY}=\textrm{Cov}(X,Y)\)
\item
  However, in reality, the quantities \(\sigma_{X}^{2}\),
  \(\sigma_{Y}^{2}\) and \(\sigma_{XY}\) are unknown, and need to be
  estimated
\end{itemize}

\subsection{Illustration I: estimate}\label{illustration-i-estimate}

\begin{itemize}
\tightlist
\item
  With estimates \(\hat{\sigma}_{X}^{2}\), \(\hat{\sigma}_{Y}^{2}\) and
  \(\hat{\sigma}_{XY}\) for \(\sigma_{X}^{2}\), \(\sigma_{Y}^{2}\) and
  \(\sigma_{XY}\), respectively, we have the \emph{plug-in estimate}
  \[\hat{\alpha}=\frac{\hat{\sigma}_{Y}^{2}-\hat{\sigma}_{XY}}{\hat{\sigma}_{X}^{2}+\hat{\sigma}_{Y}^{2}-2\hat{\sigma}_{XY}}\]
  for the \emph{optimal but unknown} solution \[
  \alpha=\frac{\sigma_{Y}^{2}-\sigma_{XY}}{\sigma_{X}^{2}+\sigma_{Y}^{2}-2\sigma_{XY}}
  \]
\item
  How accurate is \(\hat{\alpha}\)? Can we estimate the standard error
  of \(\hat{\alpha}\)?
\end{itemize}

\subsection{Illustration I: estimate}\label{illustration-i-estimate-1}

\begin{itemize}
\item
  If \(\hat{\sigma}_{X}^{2}\), \(\hat{\sigma}_{Y}^{2}\) and
  \(\hat{\sigma}_{XY}\) are accurate, then so should be \(\hat{\alpha}\)
\item
  How to assess the accuracy of \(\hat{\alpha}\) (via the accuracy of
  \(\hat{\sigma}_{X}^{2}\), \(\hat{\sigma}_{Y}^{2}\) and
  \(\hat{\sigma}_{XY}\)) if we have only a sample of size \(n\) at hand?
\item
  Mini discussion on the question above: Case 1 ``\(n\) small'', Case 2
  ``\(n\) moderate'', and case 3 ``\(n\) large''
\end{itemize}

\subsection{Illustration I: simulated
samples}\label{illustration-i-simulated-samples}

If we know the population distribution, we can simulate samples:

\begin{center}\includegraphics[width=0.8\linewidth]{fig5_9} \end{center}

\subsection{Illustration I: simulated
samples}\label{illustration-i-simulated-samples-1}

\begin{itemize}
\item
  Suppose we simulate \(B=1000\) \emph{independent} samples for
  \((X,Y)\) (if we knew the truth), we will have \(B\) estimates
  \(\hat{\alpha}_{j},j=1,\ldots,B\) of \(\alpha\)
\item
  The sample mean
  \(\bar{\alpha}=\frac{1}{B}\sum_{j=1}^{B}\hat{\alpha}_{j}\) (of
  \(\hat{\alpha}_{j}\)'s) should be close to \(\alpha\)
\item
  The sample standard deviation
  \[s\left(  \hat{\alpha}\right)  =\sqrt{\frac{1}{B-1}\sum_{j=1}^{B}\left(
  \hat{\alpha}_{j}-\bar{\alpha}\right)  ^{2}}\] (of of
  \(\hat{\alpha}_{j}\)'s) should be close to
  \(\sigma_{\hat{\alpha}}=\sqrt{\textrm{Var}(\hat{\alpha)}}\)
\end{itemize}

\subsection{Illustration I: truth and
estimate}\label{illustration-i-truth-and-estimate}

\begin{itemize}
\item
  Truth: \(\sigma_{X}^{2}=1\), \(\sigma_{Y}^{2}=1.25\),
  \(\sigma_{XY}=0.5\) and \(\alpha=0.6\)
\item
  Estimates based on \(B=1000\) simulated, independent samples:
  \(\bar{\alpha}=0.5996\) and \(s\left( \hat{\alpha}\right)=0.083\)
\item
  Interpretation: for a random sample from the population, we would
  expect \(\hat{\alpha}\) to differ from \(\alpha\) by approximately
  \(0.08\) on average
\end{itemize}

\emph{Note:} is the ``1 standard deviation'' rule sensible?

\section{Bootstrap: definition and
applications}\label{bootstrap-definition-and-applications}

\subsection{Simulation and
double-dipping}\label{simulation-and-double-dipping}

\begin{itemize}
\tightlist
\item
  \emph{Simulated from the truth}: when we know a data generating
  process, we can simulate samples to estimate a statistic on the
  process. However, if we know the truth, why do we need to estimate the
  statistic?
\item
  \emph{Simulated from the estimate}: with a sample from a data
  generating process, we can estimate the process, use the estimated
  process to generate samples, and use the generated samples to estimate
  a statistic
\item
  \emph{Resampling from the sample}: sample randomly from a sample from
  a data generating process, regard the sampled observations as a new
  data set, and use them to estimate a statistic
\end{itemize}

\subsection{Bootstrap: definition}\label{bootstrap-definition}

In order to assess the distributional properties of an estimate of a
statistic, the bootstrap

\begin{itemize}
\tightlist
\item
  takes a subset of a given data set \emph{as if it is a set of new
  observations independent of the given data set}
\item
  uses the subset to obtain an estimate of the statistic
\item
  does so \emph{repeatedly and independently using different subests of
  the given data set}
\item
  take the \emph{empirical distribution of estimates obtained from these
  subsets as an estimate of the distribution of the estimate} of the
  statistic
\end{itemize}

\subsection{Bootstrap: procedure}\label{bootstrap-procedure}

\begin{itemize}
\tightlist
\item
  Given a sample of size \(n\), let \(\hat{\alpha}\) be an estimate of a
  statistic \(\alpha\) obtained from the sample
\item
  \emph{Sample randomly with replacement} from the sample to obtain
  \(n\) observations, and do this independently \(B\) times to obtain
  \(B\) bootstrap samples \(S_{j},j=1,\ldots,B\)
\item
  Let \(\hat{\alpha}_{j}\) be the estimate of \(\alpha\) obtained from
  \(S_{j}\). Then \emph{the empirical distribution \(G\) of
  \(\hat{\alpha}_{j},j=1,\ldots,B\) is used as the (true) distribution
  of \(\hat{\alpha}\), and statistics about \(\hat{\alpha}\) are
  obtained from \(G\)}
\end{itemize}

\subsection{Bootstrap: graphical
illustration}\label{bootstrap-graphical-illustration}

\begin{center}\includegraphics[width=0.9\linewidth]{fig5.11} \end{center}

\subsection{Bootstrap: statistics}\label{bootstrap-statistics}

\begin{itemize}
\tightlist
\item
  The \emph{(bootstrap) estimated mean} of \(\hat{\alpha}\) is
  \(\bar{\alpha}=\frac{1}{B}\sum_{j=1}^{B}\hat{\alpha}_{j}\)
\item
  The \emph{(bootstrap) estimated variance} of \(\hat{\alpha}\) is \[
  \textrm{SE}^2\left(  \hat{\alpha}\right)  = {(B-1)}^{-1}\sum_{j=1}^{B}\left(
  \hat{\alpha}_{j}-\bar{\alpha}\right)^{2}
  \]
\item
  For \(\alpha \in (0,1)\), the \emph{(bootstrap)
  \((1-\alpha)\times 100\) percent confidence interval} for
  \(\hat{\alpha}\) is \((c_L,c_U)\), where \(c_L\) is the
  \(\{0.5\alpha\times 100\}\)th percentile of \(G\), and \(c_L\) is the
  \(\{(100-0.5\alpha)\times 100\}\)th percentile of \(G\)
\end{itemize}

\subsection{Illustration of bootstrap}\label{illustration-of-bootstrap}

Bootstrap applied to a sample;
\(\textrm{SE}^2\left( \hat{\alpha}\right)=0.087\): illusion or
excellence?

\begin{center}\includegraphics[width=0.9\linewidth]{fig5_10} \end{center}

\section{Boostrapping linear
regression}\label{boostrapping-linear-regression}

\subsection{Linear regression}\label{linear-regression}

\begin{itemize}
\item
  Model:
  \(Y=\beta_0+\beta_1 X_1 + \beta_2 X_2 + \ldots + \beta_p X_p + \varepsilon\)
  with \(E(\varepsilon)=0\) and \(\textrm{Var}(\varepsilon)=\sigma^2\)
\item
  Observations: \((y_i,x_{1i},x_{2i},\ldots,x_{pi}),i=1,\ldots,n\),
  where \(x_{ji}\) is the \(i\)th observation for \(X_j\)
\item
  Estimate:
  \(\hat{y}=\hat{\beta}_0+\hat{\beta}_1 X_1 + \hat{\beta}_2 X_2 + \ldots + \hat{\beta}_p X_p\)
\item
  Fit:
  \(\hat{y}_i=\hat{\beta}_0+\hat{\beta}_1 x_{1i} + \hat{\beta}_2 x_{2i} + \ldots + \hat{\beta}_p x_{pi}+\varepsilon_i\)
\item
  Residuals: \(e_i = y_i - \hat{y}_i\)
\end{itemize}

\subsection{Bootstrapping from sample}\label{bootstrapping-from-sample}

Set \({\boldsymbol{\beta}}=({\beta}_0,{\beta}_1,\ldots,{\beta}_p)\) and
\(\hat{\boldsymbol{\beta}}=(\hat{\beta}_0,\hat{\beta}_1,\ldots,\hat{\beta}_p)\)

\begin{itemize}
\tightlist
\item
  Sample from data generating process:
  \[S=\{\mathbf{z}_i=(y_i,x_{1i},x_{2i},\ldots,x_{pi}),i=1,\ldots,n\}\]
\item
  Sample with replacement \(n\) observations from \(S\) and repeat this
  independently to obtain \(B\) subsets \(S_j,j=1,\ldots,B\)
\item
  Obtain \(\hat{\boldsymbol{\beta}}_j\) from \(S_j\) for each
  \(j=1,\ldots,B\)
\item
  Use the empirical distribution of
  \(\hat{\boldsymbol{\beta}}_j,j=1,\ldots,B\) as the distribution of
  \(\hat{\boldsymbol{\beta}}\)
\end{itemize}

\subsection{Bootstrapping residuals}\label{bootstrapping-residuals}

\begin{itemize}
\tightlist
\item
  Residuals: \(R=\{e_i=y_i-\hat{y}_i,i=1,\ldots,n\}\)
\item
  Sample with replacement \(n\) observations from \(R\) to obtain \(B\)
  sets of residuals \(R_j=\{e_i^{(j)},i=1,\ldots,n\}\)
\item
  For each \(j\), set \(y_i^{(j)}=\hat{y}_i+e_i^{(j)}\) and fit the
  model with observations
  \[S_j=\{\mathbf{z}_i^{(j)}=(y_i^{(j)},x_{1i},x_{2i},\ldots,x_{pi}),i=1,\ldots,n\}\]
  and obtain estimate \(\hat{\boldsymbol{\beta}}_j\)
\item
  Use the empirical distribution of
  \(\hat{\boldsymbol{\beta}}_j,j=1,\ldots,B\) as the distribution of
  \(\hat{\boldsymbol{\beta}}\)
\end{itemize}

\subsection{Bootstrapping samples or
residuals}\label{bootstrapping-samples-or-residuals}

\begin{itemize}
\item
  Asymptotically (and under some conditions), bootstrapping samples and
  bootstrapping residuals are equivalent
\item
  Bootstrapping samples is less sensitive to model misspecification
\item
  Bootstrapping samples may be less sensitive to the assumptions
  concerning independence or exchangeability of the error terms
\end{itemize}

\section{Boostrap: failures}\label{boostrap-failures}

\subsection{Boostrap failures}\label{boostrap-failures-1}

Bootstrap can fail

\begin{itemize}
\tightlist
\item
  when sample size is too small
\item
  for estimating extremal statistics
\item
  when observations are dependent
\item
  for survey sampling
\end{itemize}

\emph{Note:} the book ``Bootstrap methods: a guide for practitioners and
researchers'' by Michael R. Chernick contains more information on this.

\subsection{License and session
Information}\label{license-and-session-information}

\href{http://math.wsu.edu/faculty/xchen/stat412/LICENSE.html}{License}

\begin{Shaded}
\begin{Highlighting}[]
\NormalTok{>}\StringTok{ }\KeywordTok{sessionInfo}\NormalTok{()}
\NormalTok{R version }\FloatTok{3.5.0} \NormalTok{(}\DecValTok{2018-04-23}\NormalTok{)}
\NormalTok{Platform:}\StringTok{ }\NormalTok{x86_64-w64-mingw32/}\KeywordTok{x64} \NormalTok{(}\DecValTok{64}\NormalTok{-bit)}
\NormalTok{Running under:}\StringTok{ }\NormalTok{Windows }\DecValTok{10} \KeywordTok{x64} \NormalTok{(build }\DecValTok{19041}\NormalTok{)}

\NormalTok{Matrix products:}\StringTok{ }\NormalTok{default}

\NormalTok{locale:}
\NormalTok{[}\DecValTok{1}\NormalTok{] LC_COLLATE=English_United States}\FloatTok{.1252} 
\NormalTok{[}\DecValTok{2}\NormalTok{] LC_CTYPE=English_United States}\FloatTok{.1252}   
\NormalTok{[}\DecValTok{3}\NormalTok{] LC_MONETARY=English_United States}\FloatTok{.1252}
\NormalTok{[}\DecValTok{4}\NormalTok{] LC_NUMERIC=C                          }
\NormalTok{[}\DecValTok{5}\NormalTok{] LC_TIME=English_United States}\FloatTok{.1252}    

\NormalTok{attached base packages:}
\NormalTok{[}\DecValTok{1}\NormalTok{] stats     graphics  grDevices utils     datasets  methods  }
\NormalTok{[}\DecValTok{7}\NormalTok{] base     }

\NormalTok{other attached packages:}
\NormalTok{[}\DecValTok{1}\NormalTok{] knitr_1}\FloatTok{.21}

\NormalTok{loaded via a }\KeywordTok{namespace} \NormalTok{(and not attached):}
\StringTok{ }\NormalTok{[}\DecValTok{1}\NormalTok{] compiler_3}\FloatTok{.5.0}  \NormalTok{magrittr_1}\FloatTok{.5}    \NormalTok{tools_3}\FloatTok{.5.0}    
 \NormalTok{[}\DecValTok{4}\NormalTok{] htmltools_0}\FloatTok{.3.6} \NormalTok{yaml_2}\FloatTok{.2.0}      \NormalTok{Rcpp_1}\FloatTok{.0.3}     
 \NormalTok{[}\DecValTok{7}\NormalTok{] stringi_1}\FloatTok{.2.4}   \NormalTok{rmarkdown_1}\FloatTok{.11}  \NormalTok{stringr_1}\FloatTok{.3.1}  
\NormalTok{[}\DecValTok{10}\NormalTok{] xfun_0}\FloatTok{.4}        \NormalTok{digest_0}\FloatTok{.6.18}   \NormalTok{evaluate_0}\FloatTok{.12}  
\end{Highlighting}
\end{Shaded}


\end{document}
